\documentclass[swedish]{maucsthesis}

\usepackage[utf8]{inputenc}
\usepackage{lipsum}
\usepackage{nomencl}
\usepackage{xcolor}
%\usepackage[parfill]{parskip}

\begin{document}

\swedishtitle{Lokförarnas kommunikations- och\\[10pt] informationsbehov inom C-DAS}
\englishtitle{Information and communication needs for train drivers within SFERA}

\author{Anthon Haväng}

\level{Grundnivå}
\credits{15}
\degree{Kandidatexamen, 180 hp}
\subject{Datavetenskap}
\studyprogram{Systemutvecklare}

\seminardate{2024-05-27}

\supervisor{Åse Jevinger}
\examiner{Jesper Larsson}
\maketitle %

\begin{sammanfattning}
\lipsum[1-2]
\end{sammanfattning}

\begin{abstract}
\lipsum[3-4]
\end{abstract}

% This is sort of magic, do not touch
\ifodd\value{page}\else\mbox{}\newpage\fi
\tableofcontents
\newpage
\startpagecount

\section{Begrepp och förkortningar}
\input{a_begrepp_och_forkortningar/a_begrepp_och_forkortningar}

\section{Inledning}

\subsection{Bakgrund}
Trafikverket vill optimera sin operativa trafikledning för att bättre kunna hantera störningar, förbättra kapacitet och minska slitage på infrastrukturen ~\cite{Angel2020}, ~\cite{Joborn2023}.
En nyckelfaktor i detta är att möjliggöra för trafikledningscentralernas tågklarerare att i realtid kunna planera om ett enskilt tågs körplan~\cite{Joborn2023}.
För att realisera detta krävs omfattande förändringar i Trafikverkets operativa ledningssystem, men även i järnvägsföretagens tekniska utrustning ombord på tågen~\cite{Joborn2023}.

I dagens samhälle skrivs bakgrund i presens då en nulägessituation visas, ~\cite{Lindmark2022}, ~\cite{Angel2020}, ~\cite{Olsson2023}.


\subsection{Problemformulering}
\paragraph{}
\hspace-11ptSäkerheten är alltid av högsta prioritering för en lokförare; hen måste upprätthålla ett koncentrerat fokus framåt under körning.
Förutom att fokusera på banan framåt måste även tågets framförande (pådrag, rullning, broms etc.), status på ombordsystem, diverse inkommande och utgående samtal och meddelanden (från järnvägsföretagets trafikledning, Trafikverkets trafikledning samt tjänstgörande ombordpersonal) skänkas fokus och det är upp till föraren att prioritera och göra rätt beroende på vad situationen kräver.
När då ytterligare ett moment som har hög prioritet tillförs bland nämnda distraktioner bör det ske på ett sådant sätt att föraren \emph{vill} och \emph{kan} följa de anvisningar som förordas via C-DAS\@.
Trafikverkets pilottester har påvisat att förarens förståelse av kontexten för den nya angivna körplanen kan spela in i viljan att realisera den~\cite{Angel2020, Lindmark2022}.
I förlängningen kan detta påverka järnvägsföretagens implementeringstakt av C-DAS\@.
%Att framföra tåg är ett hantverk; e
En skicklig förare har efter flera års yrkesutövande djuplodande kännedom om fordonet som framförs och banans särskilda egenskaper och därmed vilken körning som fordras därefter [källa behövs].
En förare som förstår \emph{varför} inlärd erfarenhetsbaserad hastighetsprofil kan behöva frångås har större vilja att köra inom den mottagna anvisningen.
Då realtidsjusteringarna av tidtabeller ofta berör flera tåg kommer de som inte tar emot och följer dessa anvisningar få en stor negativ inverkan på trafikflödet~\cite{Joborn2023}.
Dels för att det inte går att se hur tåget framförs men även att det inte kan delges ny tidtabellsinformation.
Det ligger därför i Trafikverkets intresse att så många järnvägsföretag som möjligt implementerar en C-DAS-lösning.


\subsection{Frågeställning och syfte}
%Syftet med studien är att undersöka möjligheterna att utveckla kommunikation och informationsutbyte mellan lokförare med hjälp av C-DAS, bortom både det grundläggande protokoll som Trafikverket tillämpar och som i sin tur bygger på vad som fastställs enligt SFERA\@.
\paragraph{}
%\hspace{-11pt}ptptSyftet med studien är att utvärdera och vidareutveckla den av Trafikverket fastställda digitala kommunikationen och informationsutbytet mellan lokförare och tågklarerare samt undersöka eventuella förbättringsmöjligheter avseende C-DAS-kommunikationens innehåll som ännu ej prövats.
%\hspace{-11pt}ptptSyftet med studien är att utvärdera fastställd digital kommunikation och informationsutbyte mellan lokförare och tågklarerare inom C-DAS samt undersöka eventuella förbättringsmöjligheter avseende C-DAS-kommunikationens innehåll som ännu ej prövats om det kan identifieras.
\hspace{-11pt}Syftet med studien är att undersöka i vilken utsträckning fastställd digital kommunikations- och informationsutbyte mellan lokförare och tågklarerare tillgodoser lokförares informations- och kommunikationsbehov.
\paragraph{}
\hspace{-11pt}För att den nya paradigmen med C-DAS/TMS ska lyckas behöver en så stor del som möjligt av tågtrafiken framföras med aktiv C-DAS\@.
Det ligger därför i Trafikverkets intresse att järnvägsföretagen tillämpar en C-DAS-lösning i sin verksamhet.
%Om datautbytet mellan lokförare och tågklarerare kan förbättras så att förarnas förtroende för C-DAS ökar skulle det kunna leda till en högre grad av efterlevnad av hastighetsanvisningar hos förarna.
%Om datautbytet mellan lokförare och tågklarerare kan förbättras så att förarnas förtroende för C-DAS ökar skulle det kunna leda till en högre grad av efterlevnad av hastighetsanvisningar hos förarna.
Om datautbytet mellan lokförare och tågklarerare tillämpas så att förarnas informations- och kommunikationsbehov inte tillgodoses skulle det kunna leda till lägre förtroende och en lägre grad av efterlevnad av hastighetsanvisningar hos förarna.
Omvänt skulle ett tillgodosett informationsbehov med högt förtroende hos förarna i sin tur kunna leda till högre grad av efterlevnad av hastighetsanvisningar.
I förlängningen skulle detta kunna leda till snabbare införandetakt av C-DAS hos järnvägsföretagen samt good-will för Trafikverket.
\begin{enumerate}
    \item Hur ser benägenheten ut för lokförare att följa hastighetsrekommendationer från en DAS-applikation idag?
    \item Hur väl tillgodoser befintlig standard för datautbyte inom C-DAS (SFERA) lokförares informationsbehov avseende efterlevnad av hastighetsanvisningar?
%   \item Hur ser informationsbehovet ut för lokförare gällande följandet av C-DAS-anvisningar jämfört med det som erbjuds enligt SFERA-standarden?
%   \item Tillgodoses detta behov av SFERA-standarden idag?
    \item Hur kan kommunikation- och informationsutbyte mellan lokförare och tågklarerare förbättras genom att gå utöver det grundläggande protokoll Trafikverket tillämpar idag (SFERA)?
%   \item Vilka data från lokförare till TMS som inte täcks in i dagens informationskedja kan stärka tågklarerarnas beslutsstöd och i förlängningen operativa förmåga?
    \item Hur kan gamification användas för att stödja och motivera lokförare i deras användande av C-DAS-verktyg?
\end{enumerate}

%\subsection{Hypotes}
%\begin{enumerate}
    \item Benägenheten för lokförare att följa hastighetsrekommendationer från en C-DAS-klient riskerar att bli låg om tillämpningen av C-DAS-klienten görs oaktsamt.
    \item Utvecklare av en C-DAS-lösning bör beakta lokförarens perspektiv och tillämpa en applikation som inte distraherar föraren i onödan och som upprätthåller ett högt förtroendekapital hos lokförarkåren.
    \item Hastighetsrekommendationer bör återges först när tågfärden efter start har nått den övre hastighetsgränsen för sin journey profile och med siffror ange inom vilket spann/hastighetsfönster tåget kan röra sig (snarare än en exakt hastighet) eller om coasting (frirull) ska tillämpas.
    \item Hastighetsrekommendationer bör återges först när tågfärden efter start har nått den övre hastighetsgränsen för sin journey profile och med siffror ange inom vilket spann/hastighetsfönster tåget kan röra sig (snarare än en exakt hastighet) eller om coasting (frirull) ska tillämpas.
    \item För förare av godståg kan det vara svårare att uppfylla tågfärdens journey profile och behovet av frihet och manöverutrymme gällande tågets framförande är således större.
\end{enumerate}

\section{Litteraturstudie}

\subsection{Urval}
\paragraph{}
Ett antal urvalskriterier togs fram för att nå ett relevant resultat från litteraturstudien.
Valet av databaser grundade sig i att det fanns relativt lite publicerad forskning kopplat till uppsatsens område.
Exempelvis saknade ACM:s databas relevant publicerad forskning och användes därför inte.
Sökträffarna från IEEE:s databas täckte forskningsområdet och har därför använts.
Sökning i Google scholar erhöll många sökträffar vilket utgjorde ett komplement till resultaten från IEEE:s databas.
Två söksträngar användes: \emph{(\lq driver advisory system\rq AND railway)} och \emph{(\lq driver advisory systems\rq AND railway)}.
\begin{enumerate}
    \item Inklusionskriterium
    \begin{enumerate}
        \item Skriven på svenska eller engelska
        \item Publicerad i vedertaget vetenskapligt forum
        \item Beskriver SFERA-standarden eller annat kommunikationsprotokoll mellan lokförare och tågklarerare
        \item Beskriver datat som tillhandahålles av TMS
    \end{enumerate}
    \item Exklusionskriterium
    \begin{enumerate}
        \item Publicerad innan 2010
        \item Berör utomeuropeisk järnväg
    \end{enumerate}
\end{enumerate}
\paragraph{}
Urvalsprocessen innehöll fyra steg.
Första steget innebar att utesluta verk baserat på titel.
Här sorterades eventuella dubbletter från annan databas bort.
För andra steget lästes abstract och för tredje steget lästes utöver abstract sammanfattning/conclusion (motsv.).
Kvarvarande artiklar lästes i sin helhet.
För samtliga steg prövades artiklarna mot inklusions- och exklusionskriterierna för eventuell uteslutning.
Antal sökträffar per steg redovisas nedan i tabell~\ref{tab:UrvalsprocessIEEE} för IEEE:s databas samt tabell~\ref{tab:UrvalsprocessGoogleScholar} för Google scholar:s databas.
Utöver kvarvarande artiklar som återstod efter litteraturstudiens urvalsprocess ingick \emph{n} interna dokument från Trafikverket.
Sammanlagt inkluderades \emph{n} verk i litteraturstudien.
\begin{table}[hbt!]
    \centering
    \begin{tabular}{|l|c|c|}
        \hline
        Urvalssteg & Antal träffar \\
        \hline
        0. Träffar & 17 \\
        1. Titel & 12 \\
        2. Abstract & 7 \\
        3. Slutsats & 4 \\
        4. Hela artikeln & 3 \\
        \hline
    \end{tabular}
    \caption{Urvalsprocess IEEE}
    \label{tab:UrvalsprocessIEEE}
\end{table}\begin{table}[hbt!]
    \centering
    \begin{tabular}{|l|c|}
        \hline
        Urvalssteg & Antal träffar \\
        \hline
        0. Träffar & 386 \\
        1. Titel & ? \\
        2. Abstract & ? \\
        3. Slutsats & ? \\
        4. Hela artikeln & ? \\
        \hline
    \end{tabular}
    \caption{Urvalsprocess Google scholar}
    \label{tab:UrvalsprocessGoogleScholar}
\end{table}

\subsection{Forskningsläget idag}
\subsubsection{C-DAS}

CATO:
~\cite{Yang2013}


\subsubsection{Gamification}

Litteraturstudie av dessa artiklar redovisas med ingående under~\ref{subsec:dokumentanalys} som del av studiens datainsamling.

\section{Metod och etisk analys}

\subsection{Metodbeskrivning}

\subsubsection{Design science in information systems research (lämplig rubrik?)}
\paragraph{}
\hspace{-11}För att besvara forskningsfrågorna har Hevners metod för design science research tillämpats~\cite{Hevner2004}.
Metoden går ut på att på systematiskt vis skapa en innovativ IT-artefakt som effektivt löser ett tydligt business-problem~\cite{Hevner2004}.
Artefakten utvärderas och förbättras iterativt under studiens gång för att uppnå optimalt resultat~\cite{Hevner2004}.
Det cykliska förhållningssättet kan innefatta artefaktens skapande, tillämpning eller utvärdering.
Utöver att leda till förbättring av befintliga verktyg, system eller processer leder metoden även till djupare förståelse för det problemområde artefakten tar sikte på att förbättra~\cite{Hevner2004}.
Metoden har ett tydligt fokus på praktisk problemlösning~\cite{Hevner2004}.
\paragraph{}
\hspace{-11}En nyckelfaktor med Hevners design science research är vikten av att den är både relevant och rigorös~\cite{Hevner2004}.
%Med relevans avses vikten av att det finns ett verkligt problem som adresseras.
Med relevans avses vikten av att adressera ett verkligt problem som faktiskt finns någonstans, exempelvis i en organisation.
Målet ska vara att ta fram en teknisk lösning som på ett betydelsefullt och effektivt sätt löser problemet~\cite{Hevner2004}.
Med rigorositet avses stringens i tillämpningens alla delar för artefakten.
Genom att tillämpa och noggrant följa befintliga teorier och metoder säkerställs rigorositet~\cite{Hevner2004}.
Detta gäller allt från forskningen till design, planering och själva framtagandet av lösningen.
På dessa två hörnstenar vilar Hevners metod.
\paragraph{}
\hspace{-11}En uppsättning riktlinjer listas som kan följas i samband med både forskningen och framtagandet av artefakten för att underlätta processen att uppnå relevans och rigorositet~\cite{Hevner2004}.
Dessa riktlinjer är snarare vägledning än något som måste följas till punkt och pricka.
Däremot kan de vara ett stöd i att en forskningsstudie bedrivs metodiskt och systematiskt samt att IT-artefakten effektivt löser det identifierade problemet.
Vidare kan riktlinjerna hjälpa författaren med framtagandet av IT-artefakten.
\subsubsection{Urval}
\lipsum[2]
\subsubsection{Datainsamling}
\lipsum[3]

\subsection{Metoddiskussion}
\subsubsection{DS-forskning (Hevner)}
Valet av metod grundade sig i att DS-forskning vilar på en innovativ grund, och paradigmen med järnvägstrafik med C-DAS-stöd är ung och på många ställen ej ännu införd.
Innovationsfokuset skulle kunna lämpa sig väl för kontexten C-DAS då det inte är säkert att de mest optimala lösningar ännu framtagits eller täcks av befintlig teknik.
Design-science research drivs av relevans och trovärdighet vilket passar problem- och frågeställningen bra.
Relevansen fastställs i att verksamhetsproblemet är tydligt utstakat och väl förankrat i verkligheten.
Trovärdigheten täcks av att datautbytesprotokollet utvärderas, valideras och förfinas.
\subsubsection{Fallstudie (Oates)}
En fallstudie såsom definierat enligt Oates lämpar sig väl för djuplodande analys av särskilt specifika problem~\cite{Oates2005}.
Kvalitativ datainsamling i form av semi-strukturella intervjuer samt dokumentanalys såsom de tillämpats inom ramarna för denna studie skulle kunna lämpa sig väl även i en fallstudie.
Problemet är att en fallstudie i grunden tar fasta på redan existerande system eller fenomen.
C-DAS är inte fullt utrullat i många länder idag vilket skulle kunna ses som mindre gynnsamma förutsättningar för att genomföra en fallstudie.
En fallstudie skulle inte nödvändigtvis mynna ut i innovativa förslag på framtagande eller, i fallet för den här studien, förbättringspotential av det innevarande protokoll från lokförares perspektiv.
%För fallet med C-DAS är systemet inte vida etablerat varför det skulle kunna argumenteras för att fallstudie inte vore optimal forskningsmetod.

\subsection{Etisk analys}
Kan det missbrukas?
Åldersdiskriminering.
Deadline någon gång mellan 29/4 och 7/5.

\section{Resultat}
\lipsum[1]

\section{Analys}
\lipsum[2]

\section{Diskussion}
\lipsum[3]

\section{Slutsatser och vidare forskning}
\lipsum[4]

%
% Do not change
\newpage
\addcontentsline{toc}{section}{Referenser}

\bibliographystyle{plain}    % Simplest style to use
\bibliography{main}  % Reads the bibliography.bbl file generated by BibTeX

\end{document}